\section{Methodology}

The core objective of our methodology is to isolate the lighting load from a building's aggregate, hourly electricity meter data without relying on any sub-metered ground truth. We achieve this through a fully unsupervised, two-stage deep learning framework. Stage 1 focuses on feature extraction and clustering to identify distinct operational regimes. Stage 2 performs the actual signal disaggregation, separating the aggregate profile into a continuous base load (primarily HVAC) and discrete, occupancy-driven pulses (primarily lighting).

\subsection{Data Preprocessing}
The foundation of our analysis is the public ASHRAE Great Energy Predictor III (GEPIII) dataset \cite{ashrae_2020}. This extensive dataset comprises three years of hourly meter readings from 1,448 non-residential buildings located across 16 distinct climate zones globally.

Our preprocessing pipeline consists of the following steps:
\begin{enumerate}
    \item \textbf{Filtering:} We isolate the data to include only electricity meters (meter type 0), discarding chilled water, steam, and hot water readings, as our focus is on electrical load disaggregation.
    \item \textbf{Cleaning:} We remove any negative meter readings (which indicate sensor errors or anomalous net-metering events) and drop records with missing values.
    \item \textbf{Structuring:} The continuous, multi-year time series for each building is pivoted into discrete, 24-hour daily load profiles. Days containing incomplete data (fewer than 24 hourly readings) are discarded to ensure uniform input dimensions for the neural networks. This process yields approximately 500,000 complete daily profiles.
    \item \textbf{Normalization:} To enable the model to learn scale-invariant operational patterns rather than building-specific magnitudes, we apply Min-Max normalization to each daily profile independently. Let $x_t \in \mathbb{R}^{24}$ be the raw daily profile. The normalized profile $\tilde{x}_t$ is computed as:
    \begin{equation}
        \tilde{x}_{t, h} = \frac{x_{t, h} - \min(x_t)}{\max(x_t) - \min(x_t)} \quad \forall h \in \{0, \dots, 23\}
    \end{equation}
    This normalization maps all profiles to the range $[0, 1]$, preserving the shape of the load curve while eliminating absolute consumption differences between large and small facilities.
\end{enumerate}

\subsection{Stage 1: Unsupervised Clustering via 1D-CAE}
Before disaggregating the signal, it is crucial to understand the underlying operational context of each day. Commercial buildings exhibit highly variable usage patterns depending on the day of the week, holidays, and specific tenant activities. To capture these regimes, we employ a 1D Convolutional Autoencoder (1D-CAE) for unsupervised feature extraction, followed by MiniBatch KMeans clustering.

\subsubsection{1D Convolutional Autoencoder (1D-CAE)}
The 1D-CAE is designed to compress the 24-hour normalized profile into a compact, robust latent representation. The encoder network consists of two convolutional layers, each followed by a ReLU activation and a max-pooling operation. The decoder mirrors this structure using transposed convolutions to reconstruct the original 24-hour sequence.

The network is trained to minimize the Mean Squared Error (MSE) between the input profile $\tilde{x}_t$ and the reconstructed profile $\hat{x}_t$:
\begin{equation}
    \mathcal{L}_{CAE} = \frac{1}{N} \sum_{i=1}^{N} ||\tilde{x}^{(i)} - \hat{x}^{(i)}||_2^2
\end{equation}
By forcing the data through a lower-dimensional bottleneck, the CAE learns to ignore high-frequency noise and capture the dominant structural features of the load curve, such as the timing of the morning ramp-up, the duration of the peak plateau, and the evening ramp-down.

\subsubsection{Latent Space Clustering}
Once the 1D-CAE is trained, we extract the latent representations for all 500,000 profiles. We then apply MiniBatch KMeans clustering to this latent manifold. Based on empirical analysis of the sum of squared distances (elbow method), we select $k=3$ clusters. These clusters correspond to distinct operational regimes:
\begin{itemize}
    \item \textbf{Cluster 0 (Scheduled Weekday):} Characterized by a sharp morning increase, a sustained high plateau during business hours, and a sharp evening decline.
    \item \textbf{Cluster 1 (Reactive/Weekend):} Characterized by a lower overall magnitude, less pronounced peaks, and higher variance, typical of weekend usage or highly reactive, occupancy-driven control.
    \item \textbf{Cluster 2 (Continuous Operation):} Characterized by a relatively flat profile with minimal diurnal variation, indicative of facilities operating 24/7 (e.g., data centers, hospitals).
\end{itemize}

\subsection{Stage 2: N-BEATS Disaggregation}
The core innovation of our framework is the adaptation of the N-BEATS architecture \cite{nbeats_2020} for unsupervised signal decomposition. Originally developed for time series forecasting, N-BEATS utilizes a deep stack of fully connected layers with forward and backward residual links. Crucially, the architecture can be configured to output interpretable components by constraining the final layers to produce specific basis functions.

We configure our N-BEATS model with two distinct stacks: a \textbf{Trend Block} and a \textbf{Seasonality Block}.
\begin{enumerate}
    \item \textbf{Trend Block:} This block is designed to capture low-frequency, slowly varying components of the signal. In the context of a commercial building's electrical load, this corresponds primarily to the continuous base load, which is dominated by HVAC systems maintaining baseline thermal conditions, as well as always-on equipment like servers and refrigeration.
    \item \textbf{Seasonality Block:} This block is designed to capture high-frequency, periodic, and discrete components. This aligns perfectly with the operational profile of lighting and personal electronics, which exhibit sharp, occupancy-driven pulses (turning on in the morning and off in the evening).
\end{enumerate}

The model takes the 24-hour normalized aggregate profile $\tilde{x}_t$ as input. The Trend Block first processes the input, outputting a trend component $\hat{y}_{trend}$. A backward residual link computes the remaining signal:
\begin{equation}
    r_{trend} = \tilde{x}_t - \hat{y}_{trend}
\end{equation}
The Seasonality Block then processes this residual, outputting the seasonal component $\hat{y}_{seasonality}$. The final reconstruction of the aggregate signal is the sum of these two components:
\begin{equation}
    \hat{x}_{NBEATS} = \hat{y}_{trend} + \hat{y}_{seasonality}
\end{equation}

The entire network is trained end-to-end to minimize the reconstruction MSE between $\tilde{x}_t$ and $\hat{x}_{NBEATS}$. Because the architecture explicitly forces the signal through these two distinct basis expansions, the network naturally learns to separate the low-frequency base load ($\hat{y}_{trend}$) from the high-frequency lighting pulses ($\hat{y}_{seasonality}$) without requiring any sub-metered ground truth. The isolated lighting load for auditing purposes is therefore directly given by $\hat{y}_{seasonality}$.
